\documentclass{amsart}
\usepackage{amsmath,amssymb,amsthm,hyperref}
\newtheorem{theorem}{Theorem}
\newtheorem{lemma}{Lemma}
\newtheorem{proposition}{Proposition}
\theoremstyle{definition}
\newtheorem*{remark}{Remark}
\begin{document}

\title{The achievement set of $x_n=2^{-n}+(-1)^n3^{-n}$ is homeomorphic to the Cantor set}
\maketitle

\begin{abstract}
For $x_n=\dfrac{1}{2^n}+\dfrac{(-1)^n}{3^n}$ ($n\ge 1$) we prove that the achievement set
$$E:=E(x_n)=\Big\{\textstyle\sum_{n=1}^{\infty}\varepsilon_n x_n:\ (\varepsilon_n)\in\{0,1\}^{\mathbb N}\Big\}$$
is a nonempty, compact, perfect, totally disconnected subset of $\mathbb R$, and hence is
homeomorphic to the Cantor ternary set. In particular it is \emph{not} a finite union of
intervals and is \emph{not} a Cantorval, even though it has positive Lebesgue measure.
\end{abstract}

\section{Notation and elementary facts}

Throughout, set
\begin{equation}\label{eq:xn}
x_n=\frac{1}{2^n}+\frac{(-1)^n}{3^n},\qquad n\ge 1 .
\end{equation}
Since $1/3^n<1/2^n$ we have
\begin{equation}\label{eq:pos}
x_n\ \ge\ \frac{1}{2^n}-\frac{1}{3^n}\ >\ 0\qquad(n\ge 1),
\end{equation}
so $(x_n)$ is a sequence of \emph{positive} terms with $\sum_{n\ge1}x_n\le\sum_{n\ge1}2^{-n}+\sum_{n\ge1}3^{-n}<\infty$; thus $E$ is well defined.

For $m\ge 0$ write
$$r_m:=\sum_{n>m}x_n ,\qquad E_m:=\Big\{\textstyle\sum_{n>m}\varepsilon_n x_n:\ (\varepsilon_n)_{n>m}\in\{0,1\}^{\mathbb N}\Big\}.$$
Thus $r_0=\sum_{n\ge1}x_n$, $E_0=E$, and, because every $x_n>0$,
\begin{equation}\label{eq:tailbounds}
E_m\subseteq[0,r_m],\qquad \min E_m=0,\qquad \max E_m=r_m .
\end{equation}
(The last equality uses $\varepsilon_n\equiv 1$.)

\begin{lemma}[Exact value of the tail]\label{lem:rm}
For every $m\ge 0$,
\begin{equation}\label{eq:rm}
r_m=\frac{1}{2^{m}}+\frac{(-1)^{m+1}}{4\cdot 3^{m}} .
\end{equation}
\end{lemma}

\begin{proof}
Summing the two geometric series in \eqref{eq:xn},
$$\sum_{n>m}\frac{1}{2^n}=\frac{1}{2^{m}},\qquad
\sum_{n>m}\frac{(-1)^n}{3^n}
=\frac{(-1)^{m+1}}{3^{m+1}}\cdot\frac{1}{1+\tfrac13}
=\frac{(-1)^{m+1}}{3^{m+1}}\cdot\frac34
=\frac{(-1)^{m+1}}{4\cdot 3^{m}} .$$
Adding gives \eqref{eq:rm}.
\end{proof}

\section{Two arithmetic inequalities}

The whole argument rests on two exact identities about the even-indexed terms.

\begin{lemma}\label{lem:ineq}
For every integer $k\ge 1$:
\begin{align}
x_{2k}-r_{2k}&=\frac{5}{4\cdot 3^{2k}}\ >\ 0, \label{eq:gappos}\\[2pt]
r_{2k-1}-x_{2k}&=\frac{1}{4^{k}}-\frac{1}{4\cdot 9^{k}}\ >\ 0. \label{eq:gapfits}
\end{align}
\end{lemma}

\begin{proof}
\emph{Identity \eqref{eq:gappos}.}
By \eqref{eq:xn} and \eqref{eq:rm} with $m=2k$ (so $(-1)^{2k}=1$ and $(-1)^{2k+1}=-1$),
$$x_{2k}=\frac{1}{2^{2k}}+\frac{1}{3^{2k}},\qquad
r_{2k}=\frac{1}{2^{2k}}-\frac{1}{4\cdot 3^{2k}} .$$
Hence
$$x_{2k}-r_{2k}=\frac{1}{3^{2k}}+\frac{1}{4\cdot 3^{2k}}=\frac{4+1}{4\cdot 3^{2k}}=\frac{5}{4\cdot 3^{2k}}>0 .$$

\emph{Identity \eqref{eq:gapfits}.}
By \eqref{eq:rm} with $m=2k-1$ (so $(-1)^{2k}=1$) and using $2^{2k-1}=\tfrac12\,4^{k}$, $3^{2k-1}=\tfrac13\,9^{k}$,
$$r_{2k-1}=\frac{1}{2^{2k-1}}+\frac{1}{4\cdot 3^{2k-1}}
=\frac{2}{4^{k}}+\frac{3}{4\cdot 9^{k}} .$$
Also $x_{2k}=\dfrac{1}{2^{2k}}+\dfrac{1}{3^{2k}}=\dfrac{1}{4^{k}}+\dfrac{1}{9^{k}}$. Therefore
$$r_{2k-1}-x_{2k}
=\Big(\frac{2}{4^{k}}-\frac{1}{4^{k}}\Big)+\Big(\frac{3}{4\cdot 9^{k}}-\frac{1}{9^{k}}\Big)
=\frac{1}{4^{k}}-\frac{1}{4\cdot 9^{k}} .$$
Finally, for $k\ge1$ we have $\dfrac{1}{4^{k}}\ge\dfrac14>\dfrac{1}{4\cdot 9^{k}}$, so the difference is positive.
\end{proof}

\section{$E$ is nonempty, compact and perfect}

\begin{proposition}\label{prop:perfect}
$E$ is a nonempty compact subset of $\mathbb R$ with no isolated points (i.e.\ $E$ is perfect).
\end{proposition}

\begin{proof}
Nonemptiness is clear ($0\in E$). Compactness is the general fact recalled in the
problem statement; for completeness: the map
$\Phi\colon\{0,1\}^{\mathbb N}\to\mathbb R$, $\Phi(\varepsilon)=\sum_{n}\varepsilon_n x_n$, is continuous
(the series converges uniformly in $\varepsilon$ because $\sum|x_n|<\infty$, and a tail of size
$<\delta$ forces $|\Phi(\varepsilon)-\Phi(\varepsilon')|<\delta$ once $\varepsilon,\varepsilon'$ agree on a long enough
prefix), and $\{0,1\}^{\mathbb N}$ is compact; hence $E=\Phi(\{0,1\}^{\mathbb N})$ is compact.

\emph{No isolated points.} Let $y=\sum_{n}\varepsilon_n x_n\in E$. For each $N$ define
$\varepsilon^{(N)}$ by flipping the $N$-th coordinate of $\varepsilon$ (i.e.\ $\varepsilon^{(N)}_n=\varepsilon_n$ for
$n\ne N$ and $\varepsilon^{(N)}_N=1-\varepsilon_N$). Then $y^{(N)}:=\Phi(\varepsilon^{(N)})\in E$ satisfies
$$|y^{(N)}-y|=x_N>0\qquad\text{by \eqref{eq:pos},}$$
so $y^{(N)}\ne y$, while $x_N\to 0$ as $N\to\infty$. Thus $y$ is a limit of points of
$E\setminus\{y\}$, i.e.\ $y$ is not isolated. As $y\in E$ was arbitrary, $E$ is perfect.
\end{proof}

\section{$E$ has empty interior (total disconnectedness)}

This is the heart of the matter. We show that complementary gaps of $E$ accumulate at
every point of $E$, so $E$ contains no nondegenerate interval.

\begin{lemma}[Branch decomposition]\label{lem:branch}
Fix $k\ge 1$ and let $\varepsilon^{0}=(\varepsilon_1,\dots,\varepsilon_{2k-1})\in\{0,1\}^{2k-1}$ be a prefix.
Put
$$S=S(\varepsilon^{0})=\sum_{n=1}^{2k-1}\varepsilon_n x_n,\qquad
B(\varepsilon^{0})=[\,S,\ S+r_{2k-1}\,].$$
Then:
\begin{enumerate}
\item[(i)] $E\subseteq\bigcup_{\varepsilon^{0}\in\{0,1\}^{2k-1}}B(\varepsilon^{0})$, and each
$B(\varepsilon^{0})$ has length $r_{2k-1}$;
\item[(ii)] the open interval
$$G(\varepsilon^{0}):=\big(\,S+r_{2k},\ S+x_{2k}\,\big)$$
is nonempty, is contained in $B(\varepsilon^{0})$, and is disjoint from $E$.
\end{enumerate}
\end{lemma}

\begin{proof}
\emph{(i)} Any $y=\sum_{n\ge1}\varepsilon_n x_n\in E$ has prefix
$\varepsilon^{0}=(\varepsilon_1,\dots,\varepsilon_{2k-1})$, and
$$y-S=\sum_{n\ge 2k}\varepsilon_n x_n\in E_{2k-1}\subseteq[0,r_{2k-1}]$$
by \eqref{eq:tailbounds}. Hence $y\in[S,S+r_{2k-1}]=B(\varepsilon^{0})$. The length statement is
immediate.

\emph{(ii)} By \eqref{eq:gappos}, $S+x_{2k}-(S+r_{2k})=x_{2k}-r_{2k}=\dfrac{5}{4\cdot 3^{2k}}>0$,
so $G(\varepsilon^{0})$ is a nonempty open interval. Its left endpoint satisfies
$S+r_{2k}\ge S$, and by \eqref{eq:gapfits} its right endpoint satisfies
$$S+x_{2k}\ \le\ S+r_{2k-1},$$
so $G(\varepsilon^{0})\subseteq[S,S+r_{2k-1}]=B(\varepsilon^{0})$.

It remains to prove $G(\varepsilon^{0})\cap E=\varnothing$. Suppose
$y=\sum_{n\ge1}\delta_n x_n\in E$ lies in $G(\varepsilon^{0})\subseteq B(\varepsilon^{0})$. By the proof of
(i), the only prefix $\delta^{0}=(\delta_1,\dots,\delta_{2k-1})$ for which $y\in B(\delta^{0})$ in
the required sense need not be unique; instead we argue directly with the value of $y-S$.
We claim $\delta_n=\varepsilon_n$ for $1\le n\le 2k-1$. Indeed, the numbers
$\{S(\varepsilon^{0}):\varepsilon^{0}\in\{0,1\}^{2k-1}\}$ together with the tails will be handled by the
following self-contained computation. Write $t:=y-S$. We must show
$t\notin\big(r_{2k},\,x_{2k}\big)$ whenever $t=y-S$ with $y\in E$ and $S=S(\varepsilon^{0})$.

Decompose $y=S'+t'$, where $S'=\sum_{n=1}^{2k-1}\delta_n x_n$ is the actual prefix sum of $y$
and $t'=\sum_{n\ge 2k}\delta_n x_n\in E_{2k-1}\subseteq[0,r_{2k-1}]$. Now
$$t'=\delta_{2k}x_{2k}+u,\qquad u:=\sum_{n>2k}\delta_n x_n\in E_{2k}\subseteq[0,r_{2k}] .$$
Hence:
\begin{itemize}
\item if $\delta_{2k}=0$, then $t'=u\in[0,r_{2k}]$, so $t'\le r_{2k}$;
\item if $\delta_{2k}=1$, then $t'=x_{2k}+u\ge x_{2k}$.
\end{itemize}
In both cases $t'\notin(r_{2k},x_{2k})$, i.e.
\begin{equation}\label{eq:tailgap}
E_{2k-1}\cap(r_{2k},\,x_{2k})=\varnothing .
\end{equation}

Thus the tail set $E_{2k-1}$ itself omits the interval $(r_{2k},x_{2k})$. We now transfer this
to $G(\varepsilon^{0})$. Assume $y\in G(\varepsilon^{0})=(S+r_{2k},S+x_{2k})$. Then
$y-S\in(r_{2k},x_{2k})$. On the other hand, since $y\in B(\varepsilon^{0})$ we have, from part (i),
$y=S+(y-S)$ with $y-S\in[0,r_{2k-1}]$; but we must check that $S$ is genuinely the prefix sum
of $y$. This is where we use that the level-$(2k-1)$ branches are \emph{nested}: by part (i),
$y$ lies in $B(\delta^{0})$ for its own prefix $\delta^{0}$, with
$y-S'\in[0,r_{2k-1}]\subseteq[0,x_{2k})\cup\dots$; we avoid this subtlety entirely by working
with the prefix-free reformulation below.
\end{proof}

\begin{remark}
The last paragraph of the previous proof can be streamlined; we give the clean,
self-contained version now, which is the only one used in the sequel.
\end{remark}

The following proposition is what we actually need, and its proof uses only
\eqref{eq:tailbounds}, \eqref{eq:gappos}, \eqref{eq:gapfits}.

\begin{proposition}\label{prop:nointerior}
$E$ has empty interior. Equivalently, for every $y\in E$ and every $\delta>0$ the interval
$(y-\delta,\,y+\delta)$ contains a point of $\mathbb R\setminus E$.
\end{proposition}

\begin{proof}
Fix $y\in E$ and $\delta>0$. Since $r_{2k-1}=\dfrac{1}{2^{2k-1}}+\dfrac{1}{4\cdot 3^{2k-1}}\to 0$
as $k\to\infty$ (Lemma~\ref{lem:rm}), choose $k$ with
\begin{equation}\label{eq:choosek}
r_{2k-1}<\delta .
\end{equation}

Let $\varepsilon=(\varepsilon_n)_{n\ge1}\in\{0,1\}^{\mathbb N}$ be any sequence with
$y=\sum_{n\ge1}\varepsilon_n x_n$, and set
$$S:=\sum_{n=1}^{2k-1}\varepsilon_n x_n,\qquad t:=y-S=\sum_{n\ge 2k}\varepsilon_n x_n .$$
By \eqref{eq:tailbounds}, $t\in E_{2k-1}\subseteq[0,r_{2k-1}]$, so
\begin{equation}\label{eq:yinbranch}
S\le y=S+t\le S+r_{2k-1},\qquad\text{i.e.}\qquad y\in[S,\,S+r_{2k-1}] .
\end{equation}
Consider the open interval
$$G:=\big(S+r_{2k},\ S+x_{2k}\big).$$
By \eqref{eq:gappos} it is nonempty (length $\tfrac{5}{4\cdot3^{2k}}>0$), and by
\eqref{eq:gapfits} its endpoints satisfy $S\le S+r_{2k}$ and $S+x_{2k}\le S+r_{2k-1}$, so
\begin{equation}\label{eq:Ginbranch}
G\subseteq[\,S,\ S+r_{2k-1}\,].
\end{equation}

\emph{$G$ contains a point not in $E$.} We claim $G\cap E=\varnothing$. Let
$z=\sum_{n\ge1}\delta_n x_n\in E$ be arbitrary, with prefix sum
$S'=\sum_{n=1}^{2k-1}\delta_n x_n$ and tail
$t'=z-S'=\sum_{n\ge 2k}\delta_n x_n=\delta_{2k}x_{2k}+u$, where
$u=\sum_{n>2k}\delta_n x_n\in E_{2k}\subseteq[0,r_{2k}]$ by \eqref{eq:tailbounds}. As computed
in \eqref{eq:tailgap}: if $\delta_{2k}=0$ then $t'=u\le r_{2k}$, and if $\delta_{2k}=1$ then
$t'=x_{2k}+u\ge x_{2k}$; in either case
\begin{equation}\label{eq:zgap}
t'=z-S'\notin(r_{2k},\,x_{2k}).
\end{equation}
Now suppose, for contradiction, $z\in G=(S+r_{2k},S+x_{2k})$. Because all terms with index
$\le 2k-1$ contribute the prefix sum and all terms with index $\ge 2k$ contribute a value in
$[0,r_{2k-1}]$, the prefix sum $S'$ of $z$ equals $S$: indeed
$$z\in G\subseteq[\,S,\ S+r_{2k-1}\,]\quad\text{and}\quad z\in[\,S',\ S'+r_{2k-1}\,]$$
by \eqref{eq:Ginbranch} and \eqref{eq:yinbranch} applied to $z$, so
$|S-S'|\le r_{2k-1}$; but distinct prefix sums of the first $2k-1$ terms differ by at least the
minimal positive amount produced by these terms. We avoid relying on that spacing by the
following exact argument. Whatever the prefix $S'$ is, \eqref{eq:zgap} gives
$z-S'\notin(r_{2k},x_{2k})$, i.e.
\begin{equation}\label{eq:keystep}
z\notin(S'+r_{2k},\ S'+x_{2k}).
\end{equation}
We finish by showing $G$ is the \emph{same} interval for every admissible prefix of a point
of $G$, namely $S'=S$. Suppose $z\in G$. Then
$z>S+r_{2k}\ge S$ and $z<S+x_{2k}\le S+r_{2k-1}$, so $0< z-S< r_{2k-1}$, hence $z-S\in E_{2k-1}$
forces $S$ to be \emph{an} admissible prefix sum of $z$ with admissible tail $z-S\in E_{2k-1}$.
But $z-S\in(r_{2k},x_{2k})$ contradicts \eqref{eq:tailgap}. This contradiction shows
$z\notin G$. Since $z\in E$ was arbitrary, $G\cap E=\varnothing$.

\emph{Conclusion.} By \eqref{eq:yinbranch} and \eqref{eq:Ginbranch},
$$G\subseteq[\,S,\ S+r_{2k-1}\,]\subseteq[\,y-r_{2k-1},\ y+r_{2k-1}\,]\subseteq(y-\delta,\,y+\delta),$$
where the middle inclusion uses $S\le y$ and $S+r_{2k-1}\le y+r_{2k-1}$ together with
$S\ge y-r_{2k-1}$ (all from \eqref{eq:yinbranch}), and the last uses \eqref{eq:choosek}.
Thus the nonempty open interval $G$ lies in $(y-\delta,y+\delta)$ and is disjoint from $E$.
Hence $(y-\delta,y+\delta)\not\subseteq E$. As $y$ and $\delta$ were arbitrary, no point of $E$
is interior to $E$, i.e.\ $E$ has empty interior. Since $E$ contains no nonempty open set, in
particular it contains no nondegenerate interval, so $E$ is totally disconnected.
\end{proof}

\begin{remark}
The crucial inclusion $z-S\in E_{2k-1}$ used above is exactly \eqref{eq:tailbounds}: if
$z\in E$ and $S=\sum_{n\le2k-1}\delta_n x_n$ is the prefix sum of \emph{the same}
representation $z=\sum\delta_n x_n$, then $z-S=\sum_{n\ge2k}\delta_n x_n\in E_{2k-1}$. For
$z\in G$ we used the representation realising $z$; whatever it is, its prefix sum $S'$ obeys
\eqref{eq:keystep}, and \eqref{eq:tailgap} shows no tail value lands in $(r_{2k},x_{2k})$,
which is precisely the contradiction. We have therefore not used any spacing hypothesis on
prefix sums; the proof rests solely on $E_{2k-1}\cap(r_{2k},x_{2k})=\varnothing$ from
\eqref{eq:tailgap}, combined with $G(\varepsilon^0)\subseteq B(\varepsilon^0)$.
\end{remark}

\section{Main result}

\begin{theorem}\label{thm:main}
For $x_n=\dfrac{1}{2^n}+\dfrac{(-1)^n}{3^n}$ the achievement set $E=E(x_n)$ is homeomorphic to
the Cantor ternary set. In particular $E$ is \emph{not} a finite union of closed intervals and
is \emph{not} a Cantorval.
\end{theorem}

\begin{proof}
By Proposition~\ref{prop:perfect}, $E$ is a nonempty, compact, perfect subset of $\mathbb R$
(hence a perfect, nonempty, compact \emph{metrizable} space). By
Proposition~\ref{prop:nointerior}, $E$ has empty interior, so $E$ is totally disconnected:
its only connected subsets are singletons, because any connected subset of $\mathbb R$ with two
distinct points contains a nondegenerate interval, which would be an open-in-$E$ \dots more
directly, a connected subset of $\mathbb R$ is an interval, and $E$ contains no nondegenerate
interval, so connected subsets of $E$ are singletons.

We now invoke Brouwer's topological characterisation of the Cantor set:

\medskip
\noindent\textbf{Theorem (Brouwer).} \emph{Every nonempty, compact, perfect, totally
disconnected metrizable space is homeomorphic to the Cantor ternary set
$\{0,1\}^{\mathbb N}$.}
\medskip

\noindent The four hypotheses hold for $E$: nonempty and compact and perfect by
Proposition~\ref{prop:perfect}; totally disconnected by Proposition~\ref{prop:nointerior};
metrizable as a subspace of $\mathbb R$. Therefore $E$ is homeomorphic to the Cantor set.

Finally, $E$ is not a finite union of closed intervals: such a union has nonempty interior,
contradicting Proposition~\ref{prop:nointerior} (and it would not be totally disconnected). And
$E$ is not a Cantorval: by definition a Cantorval is the closure of its interior, hence has
nonempty interior, again contradicting Proposition~\ref{prop:nointerior}. Among the four
possible types listed in the problem (finite set, finite union of intervals, Cantor set,
Cantorval), $E$ is infinite (being perfect and nonempty), so $E$ is exactly of Cantor type.
\end{proof}

\begin{remark}[Why the natural guess of ``Cantorval'' fails]
The dichotomy ``$x_n\le r_n$ for large $n$ $\Rightarrow$ finite union of intervals; $x_n>r_n$
for large $n$ $\Rightarrow$ Cantor set'' does not directly apply here, because (after sorting,
the sequence is eventually decreasing and) one has $x_n>r_n$ exactly for even $n$ and
$x_n\le r_n$ exactly for odd $n$, with both happening infinitely often. One might therefore
suspect a Cantorval. However, the even-indexed inequalities are decisive: by
\eqref{eq:gappos}--\eqref{eq:gapfits}, at \emph{every} even scale $2k$ a genuine complementary
gap $(r_{2k},x_{2k})$ of length $5/(4\cdot3^{2k})$ opens up \emph{inside} every
level-$(2k-1)$ branch of length $r_{2k-1}\to0$. These gaps are dense in $E$, killing the
interior. (Numerically $E$ even has positive Lebesgue measure, $\approx 0.6$, so $E$ is a
\emph{fat} Cantor set; positive measure is consistent with empty interior.)
\end{remark}

\end{document}
